进一步，引入迹坐标 \(x=(u+u^{-1})/2\)。
在 \(|u|=1\) 的幺正切片上，\(x\in[-1,1]\)，而 \(J(u)\) 精确落在 Lee--Yang 奇异环 \((-\infty,-\tfrac14]\) 上；
该对应为“区间化压缩”与射线端点紧化的基本代数接口。

\begin{theorem}[迹--Lee--Yang 的线性分式同构]\label{thm:leyang-trace-mobius-invertibility}
设 \(J(u)=-\frac{u}{(1+u)^2}\)（\(u\in\CC\setminus\{-1\}\)）。
令
\[
x:=\frac{u+u^{-1}}{2},
\]
并假设 \(u\neq 0,-1\)（从而 \(x\) 有定义且 \(J(u)\) 有意义）。
则成立严格恒等式
\begin{equation}\label{eq:leyang-trace-mobius}
\boxed{\;
J(u)=-\frac{1}{2(1+x)}\; }.
\end{equation}
因此 \(J\) 在迹坐标 \(x\) 上退化为 Möbius 变换，其反演为
\begin{equation}\label{eq:leyang-trace-mobius-inverse}
\boxed{\;
x=-\frac{2J(u)+1}{2J(u)}\; }.
\end{equation}
特别地，当 \(|u|=1\) 且 \(u\neq -1\) 时，\(x\in[-1,1]\)，并且映射
\[
[-1,1]\ni x\longmapsto -\frac{1}{2(1+x)}\in(-\infty,-\tfrac14]
\]
为双射：端点 \(x=1\) 对应 \(J=-\tfrac14\)，而 \(x=-1\) 对应射线端点的紧化极限 \(J=-\infty\)。
\end{theorem}

\begin{proof}
由定义
\[
J(u)=-\frac{u}{(1+u)^2}.
\]
取任一平方根分支（仅用于代数化简，不引入分支依赖），有
\(
1+u=u^{1/2}(u^{1/2}+u^{-1/2})
\)，
故
\[
J(u)=-\frac{u}{u\,(u^{1/2}+u^{-1/2})^2}=-\frac{1}{(u^{1/2}+u^{-1/2})^2}.
\]
另一方面，
\[
(u^{1/2}+u^{-1/2})^2=u+u^{-1}+2=2(x+1),
\]
代入即得 \eqref{eq:leyang-trace-mobius}。
对 \eqref{eq:leyang-trace-mobius} 解出 \(x\) 即得 \eqref{eq:leyang-trace-mobius-inverse}。
当 \(|u|=1\) 且 \(u\neq -1\) 时，\(u=e^{\mathrm{i}\theta}\) 给出 \(x=\cos\theta\in[-1,1]\)，并由 \eqref{eq:leyang-trace-mobius} 得 \(J(u)\in(-\infty,-\tfrac14]\)；
线性分式在 \([-1,1]\) 上严格单调，故为双射。证毕。
\end{proof}

\endinput

